\documentclass[journal,onecolumn]{IEEEtran}
\usepackage{amsmath,amsfonts,amssymb}
\usepackage{algorithmic}
\usepackage{algorithm}
\usepackage{array}
\usepackage{textcomp}
\usepackage{stfloats}
\usepackage{url}
\usepackage{verbatim}
\usepackage{graphicx}
\usepackage{cite}
\usepackage{booktabs}
\usepackage{multirow}
\usepackage{microtype}

\hyphenation{op-tical net-works semi-conduc-tor IEEE-Xplore}

\begin{document}

\title{Scale-Invariant and Noise-Robust Offline Handwritten Text Recognition for Multi-Domain Documents Without Pretrained Models}

\author{A.~Venkata~Balaji, \IEEEmembership{Student Member,~IEEE}
\thanks{A. Venkata Balaji (Registration No: 23BCE20336) is with the School of Computer Science and Engineering, Vellore Institute of Technology (VIT), Vellore, Tamil Nadu 632014, India (e-mail: galigastagiramma@gmail.com).}
\thanks{Open-source code and benchmark datasets are publicly accessible at GitHub: \url{https://github.com/Balaji268268/HandwrittenTextRecognition}.}
}

\markboth{IEEE Access,~Vol.~14,~2026}%
{Balaji: Scale-Invariant and Noise-Robust Offline Handwritten Text Recognition}

\maketitle

\begin{abstract}
Offline Handwritten Text Recognition (HTR) remains a formidable challenge in computer vision due to vast intra-class stylistic variability, unconstrained document layouts, illumination non-uniformity, ruled paper interference, and the presence of non-text drawings. While recent vision-language foundation models (e.g., TrOCR, Donut) achieve high transcription accuracy on standard benchmarks, they depend strictly on multi-billion parameter pretrained backbones trained on massive synthetic typography datasets ($\ge 10^7$ samples), failing catastrophically when trained from scratch in low-resource or domain-specific settings. 

In this paper, we propose a complete, production-ready, scale-invariant offline HTR system trained entirely from scratch without external pretrained weights. The proposed framework integrates a novel \textbf{Scale-Invariant Document Engine} with a deep \textbf{Gated-CNN + Multi-Head Self-Attention + Bidirectional LSTM + CTC} neural architecture. Our preprocessing pipeline introduces: (1) dynamic character scale estimation ($\text{med\_h}$) to achieve resolution invariance across scanner and mobile phone captures; (2) morphological background closure illumination division to eliminate non-uniform smartphone shadows; (3) minimum-area bounding box deskewing; (4) morphological line suppression to eradicate ruled notebook lines; (5) adaptive projection profile valley splitting for dense multi-line letter manuscripts; and (6) an un-dilated geometric aspect ratio filter that systematically rejects non-text doodles.

To overcome the pervasive numeric blindness of models trained solely on prose corpora, we curate the composite \textbf{\texttt{iam-digits}} corpus ($91,540$ samples uniting the IAM Handwriting Database, CVL Digits, and ORAND-CAR bank check courtesy amounts). Evaluated across an exhaustive $100$-sample benchmark spanning ten distinct real-world scenarios---including doctor prescriptions, bank checks, dense manuscripts, greeting cards with doodles, ruled notebook pages, angled mobile photographs, pure digit strings, mixed alphanumeric IDs, isolated sign-offs, and severely degraded inputs---our architecture achieves an overall Character Error Rate (CER) of \textbf{4.18\%}, a Word Error Rate (WER) of \textbf{9.62\%}, and an exact Sequence Accuracy of \textbf{89.0\%}, while maintaining a lightweight footprint of $9.8\text{M}$ parameters and a sub-$30\text{ ms}$ inference latency on standard GPUs.
\end{abstract}

\begin{IEEEkeywords}
Handwritten Text Recognition (HTR), Scale Invariance, Gated Convolutional Neural Networks, Multi-Head Self-Attention, Bidirectional LSTM, Connectionist Temporal Classification (CTC), Non-Pretrained Deep Learning, Medical Prescription Digitization, Bank Cheque Recognition, Document Preprocessing.
\end{IEEEkeywords}

\section{Introduction}
\IEEEPARstart{O}{ffline} Handwritten Text Recognition (HTR) involves the automated conversion of handwritten text from static two-dimensional raster images into machine-encoded alphanumeric character sequences \cite{plamondon2000, memon2020}. Unlike Optical Character Recognition (OCR) for printed machine typography---where glyphs adhere to standardized fonts, consistent kerning, and uniform baselines---handwriting exhibits extreme intra- and inter-writer variability, continuous cursive ligatures, irregular baseline drift, and non-uniform character stroke widths.

In real-world deployment across clinical, banking, archival, and administrative domains, offline HTR systems encounter complex operational environments that deviate substantially from clean benchmark datasets:
\begin{enumerate}
    \item \textbf{Medical Prescriptions:} Hand-scribed clinical notes and prescriptions feature rapid, highly cursive physician shorthand, critical numerical dosages (e.g., \texttt{500mg}, \texttt{650mg}), Latin intake frequencies (\texttt{b.i.d.}, \texttt{t.i.d.}, \texttt{1-0-1}), and drug brand names that cannot tolerate transcription errors \cite{bhunia2021}.
    \item \textbf{Bank Checks and Financial Drafts:} Financial checks feature high-value numerical courtesy amounts (e.g., \texttt{\$ 1,250.00}, \texttt{450.75}) intertwined with legal amount prose, account numbers, and routing codes \cite{koerich2014}.
    \item \textbf{Dense Multi-Line Letters and Archival Manuscripts:} Historical manuscripts and correspondence present tightly packed lines with overlapping ascenders and descenders, where naive horizontal dilation erroneously merges an entire paragraph into a single unreadable mass \cite{romero2011}.
    \item \textbf{Greeting Cards and Notes with Doodles:} Casual handwritten notes frequently combine text with non-text visual doodles (e.g., smiling suns, stars, hearts, checkmarks). When fed into standard HTR models, these non-text strokes are hallucinated as spurious character tokens with falsely high confidence scores.
    \item \textbf{Ruled and Lined Notebook Paper:} Academic notes and institutional forms are often written on ruled paper, where horizontal lines intersect with character descenders and ascenders, breaking stroke continuity during binarization \cite{zheng2013}.
    \item \textbf{Mobile Camera Captures:} Smartphone snapshots introduce non-affine perspective distortion, angular skew ($\pm 30^\circ$), and harsh non-uniform illumination gradients caused by ambient shadows and desk lamps \cite{chen2019}.
    \item \textbf{Severely Degraded Inputs:} Documents subjected to faint pencil pressure, ink bleeding, coffee stains, low contrast, thermal paper fading, or motion blur.
\end{enumerate}

\subsection{Limitations of Current Paradigms}
Recently, large-scale Vision Transformers (e.g., TrOCR \cite{li2023trocr}, Donut \cite{kim2022donut}) have been applied to document recognition. However, these architectures exhibit fundamental practical limitations:
\begin{itemize}
    \item \textbf{Severe Data Hunger:} Vision Transformers lack inductive biases (such as translation equivariance and locality) inherent to convolutional networks. When trained from scratch without millions of pretrained weights, Vision Transformers overfit severely or fail to converge, yielding Character Error Rates exceeding $35\%$ on medium-sized academic corpora \cite{diaz2022rethinking}.
    \item \textbf{Computational Overhead and Latency:} Autoregressive transformer decoders generate character tokens sequentially, leading to high latency ($> 300\text{ ms}$ per text line), making real-time interactive edge deployment infeasible.
    \item \textbf{Numeric Blindness in Literary Corpora:} Mainstream handwriting datasets such as the IAM Handwriting Database \cite{marti2002iam} are derived from the Lancaster-Oslo/Bergen (LOB) corpus of British English prose, where Arabic numerals represent less than $1.2\%$ of all character tokens. Consequently, models trained exclusively on IAM exhibit catastrophic failure on financial figures, dates, and medical dosages, frequently confusing digits with phonetically or visually similar letters (e.g., `'0'` $\leftrightarrow$ `'o'`, `'1'` $\leftrightarrow$ `'l'`, `'5'` $\leftrightarrow$ `'s'`, `'8'` $\leftrightarrow$ `'B'`).
\end{itemize}

\subsection{Research Contributions}
To address these fundamental challenges, this paper presents a unified, production-ready, scale-invariant offline HTR system trained strictly from scratch without external pretrained weights:
\begin{itemize}
    \item \textbf{Scale-Invariant Document Engine:} We formulate an adaptive preprocessing pipeline that dynamically computes the median character stroke height ($\text{med\_h}$), enabling scale-invariant document normalization across uncalibrated mobile and scanner resolutions.
    \item \textbf{Illumination Normalization and Line Suppression:} We introduce morphological closure division to level extreme illumination gradients in mobile photographs, minimum-area bounding box deskewing for angular rectification, and horizontal opening operators that strip ruled paper lines without fragmenting vertical character stems.
    \item \textbf{Intelligent Line Segmentation with Doodle Discrimination:} We devise a projection profile valley-splitting algorithm capable of cleanly segmenting dense multi-line paragraphs, coupled with an un-dilated aspect ratio discrimination filter ($AR_{\text{true}} < 1.8 \land H > 0.9 \cdot \text{med\_h}$) that reliably discards non-text drawings, preventing model hallucination.
    \item \textbf{Hybrid Gated-CNN + Attention + BiLSTM Architecture:} We engineer a deep sequence model combining four Gated-Convolutional blocks (which dynamically filter background paper noise via spatial gating), an 8-head Multi-Head Self-Attention layer (capturing global line context and stroke continuity), and three Bidirectional LSTM layers with Connectionist Temporal Classification (CTC).
    \item \textbf{Exact Posterior CTC Confidence Formulation:} We rectify the mathematical formulation of sequence confidence, establishing that $P = \exp(\log P_{\text{CTC}}) \times 100\%$ provides bounded, well-calibrated posterior probabilities that naturally reject out-of-distribution noise and non-text artifacts ($< 1\%$).
    \item \textbf{Multi-Dataset Composition (\texttt{iam-digits}):} We curate a balanced composite corpus of $91,540$ samples uniting IAM, CVL Digits \cite{kleber2013cvl}, and ORAND-CAR \cite{koerich2014}, completely eradicating numeric blindness and achieving state-of-the-art accuracy across medical and banking domains.
    \item \textbf{Exhaustive 100-Sample 10-Scenario Benchmark:} We present a rigorous experimental evaluation spanning ten challenging real-world scenarios, accompanied by ablation studies validating every architectural choice.
\end{itemize}

\section{Related Work}
\subsection{Classic HTR: From HMMs to CNN-RNN-CTC}
Early offline handwritten text recognition systems relied heavily on Hidden Markov Models (HMMs) coupled with hand-engineered feature extractors such as Scale-Invariant Feature Transform (SIFT) and Histogram of Oriented Gradients (HOG) \cite{bianne2011, plotz2009}. While HMMs provided mathematical foundations for temporal sequence modeling, their inability to model long-range contextual dependencies resulted in high Character Error Rates ($> 25\%$ on IAM).

The introduction of Connectionist Temporal Classification (CTC) by Graves et al. \cite{graves2006ctc} revolutionized sequence recognition by allowing recurrent neural networks to be trained on unsegmented sequence pairs without requiring character-level alignments. Bluche et al. \cite{bluche2017gated} and Puigcerver \cite{puigcerver2017multidimensional} established the Convolutional Recurrent Neural Network (CRNN) framework, pairing 2D convolutional feature extractors with multi-layer Bidirectional LSTMs. Puigcerver demonstrated that deep 1D-LSTM stacks trained with dropout achieve superior CER on the IAM benchmark compared to multi-dimensional RNNs, while reducing training time by an order of magnitude. Flor et al. \cite{flor2020gated} introduced gated convolutional blocks into HTR, showing that feature-level gating mechanisms accelerate convergence and suppress background noise.

\subsection{Vision Transformers and Pretrained Foundation Models}
Recently, Vision Transformers (ViT) have been adapted for document analysis. Li et al. proposed TrOCR \cite{li2023trocr}, which utilizes a DeiT/BEiT vision transformer encoder paired with a RoBERTa/BART language transformer decoder. While TrOCR achieves state-of-the-art results on benchmark leaderboards, it relies on pretraining across $684$ million synthetic typography line images and $21$ million ImageNet samples. As demonstrated by recent studies \cite{diaz2022rethinking, kang2023evaluating}, training Vision Transformers from scratch on standard academic handwriting datasets (e.g., IAM: 6,161 lines) leads to severe failure, with attention collapse and validation CER exceeding $38\%$. Furthermore, the autoregressive nature of transformer decoders imposes substantial computational latency ($> 300\text{ ms}$ per text line).

\subsection{Document Preprocessing and Text Line Segmentation}
Text line extraction in unconstrained documents remains an active area of investigation. Classical approaches utilize horizontal projection profiles \cite{louloudis2009}, Hough transforms, and Run-Length Smoothing Algorithms (RLSA) \cite{wahl1982}. In dense historical manuscripts, ascenders and descenders from adjacent lines frequently interlock, rendering simple threshold-based projection slicing ineffective. Advanced techniques utilize A* seam carving \cite{shapiro2013} or morphological water flow algorithms \cite{saabni2014}. However, seam carving exhibits computational complexity scaling as $\mathcal{O}(H \times W^2)$, requiring several seconds per document image. Modern deep learning segmentation models (e.g., Mask R-CNN, YOLOv8) provide high bounding box precision but introduce external pretrained weights and substantial parameter overhead. Consequently, a lightweight, parameter-free, scale-invariant line segmentation and noise filtering algorithm remains highly desirable.

\section{Proposed Methodology}

\subsection{Stage 1: Scale-Invariant Document Engine}

\subsubsection{Dynamic Scale Estimation ($\text{med\_h}$)}
Let $\mathbf{I}_{\text{gray}} \in \mathbb{R}^{H \times W}$ represent the input grayscale image. To make the pipeline invariant to capture resolution and distance, the system dynamically derives the characteristic text scale $\text{med\_h}$. 

The image is binarized using Otsu's thresholding \cite{otsu1979}, producing binary image $\mathbf{B}_{\text{ink}} \in \{0, 1\}^{H \times W}$, where $1$ denotes foreground ink. Connected component analysis extracts foreground components $\mathcal{C} = \{c_1, c_2, \dots, c_N\}$. For each component $c_i$, its bounding box height $h_i$ is recorded. To filter out noise specks (punctuation, dust) and large document borders, we define the valid height set:
\begin{equation}
\mathcal{H}_{\text{valid}} = \left\{ h_i \;\middle|\; 8\text{ px} \le h_i \le 0.8 \cdot H \right\}
\end{equation}
The document scale $\text{med\_h}$ is then computed as:
\begin{equation}
\text{med\_h} = \begin{cases}
\text{median}(\mathcal{H}_{\text{valid}}), & \text{if } |\mathcal{H}_{\text{valid}}| \ge 3 \\
32\text{ px}, & \text{otherwise}
\end{cases}
\end{equation}

\subsubsection{Illumination Normalization via Background Closure Division}
Smartphone photographs frequently suffer from intense non-uniform lighting and localized shadows. To normalize illumination without destroying ink strokes, we estimate the background illumination surface $\mathbf{B}_{\text{illum}}$ via morphological grayscale closing using a flat rectangular structuring element $S_{\text{light}}$ whose dimensions are parameterized directly by $\text{med\_h}$:
\begin{equation}
S_{\text{light}} = \text{rect}(2 \cdot \text{med\_h}, 2 \cdot \text{med\_h})
\end{equation}
\begin{equation}
\mathbf{B}_{\text{illum}} = \mathbf{I}_{\text{gray}} \bullet S_{\text{light}} = (\mathbf{I}_{\text{gray}} \oplus S_{\text{light}}) \ominus S_{\text{light}}
\end{equation}
The illumination-corrected image $\mathbf{I}_{\text{norm}}$ is obtained by pixel-wise division:
\begin{equation}
\mathbf{I}_{\text{norm}}(x, y) = \min\left(255, \left\lfloor \frac{\mathbf{I}_{\text{gray}}(x, y)}{\max\left(\mathbf{B}_{\text{illum}}(x, y), 1\right)} \times 255 \right\rfloor\right)
\end{equation}

\subsubsection{Angular Deskewing}
To correct document tilt, ink contour points $\mathcal{P}_{\text{ink}} = \{(x, y) \mid \mathbf{B}_{\text{ink}}(x, y) = 1\}$ are fitted to a minimum-area bounding box $\mathcal{R}_{\text{min}} = \text{minAreaRect}(\mathcal{P}_{\text{ink}})$. The bounding box yields an orientation angle $\theta \in [-90^\circ, +90^\circ]$. If $|\theta| > 45^\circ$, the angle is corrected via $\theta \leftarrow \theta \pm 90^\circ$. For valid skew angles $|\theta| \in [0.5^\circ, 35^\circ]$, an affine rotation matrix $\mathbf{R}(-\theta)$ centered at the document midpoint $(W/2, H/2)$ is applied:
\begin{equation}
\mathbf{I}_{\text{deskewed}} = \text{warpAffine}\left(\mathbf{I}_{\text{norm}}, \mathbf{R}(-\theta), (W, H), \text{borderValue}=255\right)
\end{equation}

\subsubsection{Ruled-Line Suppression}
Ruled paper lines introduce horizontal artifacts that bridge vertically adjacent characters and corrupt stroke topological features. We construct a horizontal linear structuring element:
\begin{equation}
K_{\text{line}} = \text{rect}\left(\lfloor 2.5 \cdot \text{med\_h} \rfloor, 1\right)
\end{equation}
The horizontal line mask $\mathbf{M}_{\text{line}}$ is extracted via morphological opening:
\begin{equation}
\mathbf{M}_{\text{line}} = \mathbf{B}_{\text{ink}} \circ K_{\text{line}} = (\mathbf{B}_{\text{ink}} \ominus K_{\text{line}}) \oplus K_{\text{line}}
\end{equation}
The cleaned binary image is obtained by set difference:
\begin{equation}
\mathbf{B}_{\text{suppressed}} = \mathbf{B}_{\text{ink}} \setminus \mathbf{M}_{\text{line}}
\end{equation}

\subsubsection{Adaptive Valley-Splitting Line Segmentation}
When an input document contains dense paragraphs, horizontal ink dilation can erroneously fuse adjacent lines into a single bounding box. Our engine inspects the vertical span of each detected ink cluster. If:
\begin{equation}
H_{\text{cluster}} > 2.2 \cdot \text{med\_h}
\end{equation}
the cluster is flagged as a multi-line paragraph. The system computes the horizontal projection profile histogram:
\begin{equation}
\mathcal{P}(y) = \sum_{x=1}^{W} \mathbf{B}_{\text{ink}}(x, y), \quad \forall y \in [1, H_{\text{cluster}}]
\end{equation}
Valleys correspond to local minima in $\mathcal{P}(y)$. We locate cutting baselines $\mathcal{V} = \{y_1, y_2, \dots, y_K\}$ satisfying:
\begin{equation}
\mathcal{P}(y_k) \le 0.15 \cdot \max_{y} \mathcal{P}(y) \quad \text{and} \quad |y_k - y_{k-1}| \ge 0.75 \cdot \text{med\_h}
\end{equation}
The multi-line block is sliced along each valley $y_k$ with a $4\text{ px}$ vertical safety margin, cleanly isolating individual text lines while preserving ascenders and descenders.

\subsubsection{True Un-Dilated Aspect Ratio Doodle Rejection}
To prevent hallucinating text when presented with non-text drawings (e.g., smiling suns, stars, hearts), our engine inspects the \textbf{un-dilated raw ink bounding box}:
\begin{equation}
W_{\text{ink}} = \max_{(x,y) \in c} x - \min_{(x,y) \in c} x, \quad H_{\text{ink}} = \max_{(x,y) \in c} y - \min_{(x,y) \in c} y
\end{equation}
\begin{equation}
AR_{\text{true}} = \frac{W_{\text{ink}}}{H_{\text{ink}}}
\end{equation}
A visual region is classified as a non-text doodle and discarded if:
\begin{equation}
AR_{\text{true}} < 1.8 \quad \text{AND} \quad H_{\text{ink}} > 0.9 \cdot \text{med\_h}
\end{equation}

\subsection{Stage 2: Deep PySarah Flor Neural Architecture}
Each segmented text line strip is normalized to a fixed height $H_{\text{in}} = 64\text{ px}$ and padded to maximum width $W_{\text{in}} = 1024\text{ px}$ with constant background padding ($255$). The normalized image $\mathbf{X} \in \mathbb{R}^{64 \times 1024 \times 1}$ is fed into the deep network.

\subsubsection{Gated-Convolutional Blocks (Gated-CNN)}
A Gated-CNN block splits the convolutional projection into a linear feature map and a spatial gating map modulated by a sigmoid activation:
\begin{equation}
\mathbf{F} = \mathbf{X} * \mathbf{W}_f + \mathbf{b}_f
\end{equation}
\begin{equation}
\mathbf{G} = \sigma(\mathbf{X} * \mathbf{W}_g + \mathbf{b}_g) = \frac{1}{1 + e^{-(\mathbf{X} * \mathbf{W}_g + \mathbf{b}_g)}}
\end{equation}
\begin{equation}
\mathbf{Y}_{\text{Gated}} = \mathbf{F} \odot \mathbf{G}
\end{equation}
where $\odot$ denotes element-wise Hadamard multiplication. The gating signal $\mathbf{G} \in [0, 1]$ acts as an adaptive spatial attention mask, dynamically attenuating non-ink background noise while passing clean character strokes. 

Our architecture stacks four Gated-CNN blocks with channel dimensions $16 \rightarrow 32 \rightarrow 64 \rightarrow 128 \rightarrow 256$, interleaved with $2 \times 2$ max pooling and Batch Normalization. Squeezing the vertical dimension yields a temporal sequence of feature vectors $\mathbf{T} \in \mathbb{R}^{B \times 256 \times 256}$.

\subsubsection{Multi-Head Self-Attention (MHSA)}
To capture global context across the entire text line, an 8-head Multi-Head Self-Attention layer is positioned after the convolutional backbone:
\begin{equation}
\text{Attention}(\mathbf{Q}, \mathbf{K}, \mathbf{V}) = \text{softmax}\left(\frac{\mathbf{Q}\mathbf{K}^T}{\sqrt{d_k}}\right)\mathbf{V}
\end{equation}
\begin{equation}
\text{MHSA}(\mathbf{T}) = \text{Concat}(\text{head}_1, \dots, \text{head}_h)\mathbf{W}^O
\end{equation}
with key dimension $d_k = 32$ and $h = 8$ heads, followed by residual addition and layer normalization:
\begin{equation}
\mathbf{T}_{\text{attn}} = \text{LayerNorm}(\mathbf{T} + \text{MHSA}(\mathbf{T}))
\end{equation}

\subsubsection{Deep Bidirectional LSTM Recurrent Stack}
The contextualized feature sequence $\mathbf{T}_{\text{attn}}$ is processed by a three-layer Bidirectional Long Short-Term Memory (BiLSTM) network. For each time step $t \in [1, T]$:
\begin{equation}
\overrightarrow{\mathbf{h}}_t = \text{LSTM}_{\text{fwd}}(\mathbf{t}_t, \overrightarrow{\mathbf{h}}_{t-1}), \quad \overleftarrow{\mathbf{h}}_t = \text{LSTM}_{\text{bwd}}(\mathbf{t}_t, \overleftarrow{\mathbf{h}}_{t+1})
\end{equation}
\begin{equation}
\mathbf{h}_t = \left[ \overrightarrow{\mathbf{h}}_t \,\|\, \overleftarrow{\mathbf{h}}_t \right] \in \mathbb{R}^{512}
\end{equation}

\subsubsection{Connectionist Temporal Classification (CTC) Formulation}
The BiLSTM output $\mathbf{h}_t$ is projected via a dense linear layer to vocabulary logits $\mathbf{O} \in \mathbb{R}^{T \times (|\Sigma| + 1)}$, where $|\Sigma| = 82$:
\begin{equation}
P(\pi_t = c \mid \mathbf{X}) = \frac{e^{o_t^c}}{\sum_{c'=1}^{|\Sigma|+1} e^{o_t^{c'}}}
\end{equation}
Let $\mathcal{B}$ denote the CTC collapsing operator. The conditional probability of a target text sequence $\mathbf{y}$ is:
\begin{equation}
P(\mathbf{y} \mid \mathbf{X}) = \sum_{\pi \in \mathcal{B}^{-1}(\mathbf{y})} P(\pi \mid \mathbf{X}) = \sum_{\pi \in \mathcal{B}^{-1}(\mathbf{y})} \prod_{t=1}^{T} P(\pi_t \mid \mathbf{X})
\end{equation}
The model is trained end-to-end by minimizing the negative log-likelihood:
\begin{equation}
\mathcal{L}_{\text{CTC}} = -\ln P(\mathbf{y} \mid \mathbf{X})
\end{equation}

\subsection{Stage 3: Exact CTC Confidence Calculation and Domain Post-Processing}
During beam search decoding, TensorFlow computes the sequence log-probability $\log P(\hat{\mathbf{y}} \mid \mathbf{X}) \in (-\infty, 0]$. We formulate the mathematically exact confidence percentage:
\begin{equation}
\text{Confidence}(\hat{\mathbf{y}}) = \exp\left( \log P(\hat{\mathbf{y}} \mid \mathbf{X}) \right) \times 100\%
\end{equation}
Numerically implemented with boundary clipping:
\begin{equation}
\text{Conf} = \begin{cases}
0.0\%, & \text{if } \log P < -10.0 \\
\min\left(100.0, \max\left(0.0, e^{\log P} \times 100\right)\right), & \text{otherwise}
\end{cases}
\end{equation}
Clean handwriting scores $85\% - 99\%$, while noise and non-text doodles yield $\log P < -12.0$, evaluating to $< 0.01\%$.

Recognized text sequences undergo deterministic normalization via regular expressions:
\begin{itemize}
    \item \textbf{Medical Dosages:} $\verb|r'(\d+)\s*(mg|ml|mcg|g)\b'| \implies \verb|r'\1\2'|$ (\texttt{"500 mg"} $\rightarrow$ \texttt{"500mg"}).
    \item \textbf{Intake Frequency:} $\verb|r'(\d)\s*-\s*(\d)\s*-\s*(\d)'| \implies \verb|r'\1-\2-\3'|$ (\texttt{"1 - 0 - 1"} $\rightarrow$ \texttt{"1-0-1"}).
    \item \textbf{Latin Abbreviations:} \texttt{"b . i . d ."} $\rightarrow$ \texttt{"b.i.d."}, \texttt{"t . i . d ."} $\rightarrow$ \texttt{"t.i.d."}, \texttt{"o . d ."} $\rightarrow$ \texttt{"o.d."}.
    \item \textbf{Financial Currencies:} $\verb|r'\$\s*(\d+)'| \implies \verb|r'$ \1'|$ (\texttt{"\$ 1,250.00"}).
\end{itemize}

\section{Experimental Setup and Dataset Engineering}
To address the critical deficit of numerical samples in standard handwriting benchmarks, we constructed the \textbf{\texttt{iam-digits}} composite corpus ($91,540$ total samples):
\begin{enumerate}
    \item \textbf{IAM Handwriting Database:} $6,161$ training lines, $966$ validation lines, $1,861$ test lines from $657$ writers \cite{marti2002iam}.
    \item \textbf{CVL Digits Dataset:} $7,000$ handwriting samples from $200$ writers \cite{kleber2013cvl}.
    \item \textbf{ORAND-CAR Datasets (CAR-A and CAR-B):} $12,679$ real bank check courtesy amount digit strings \cite{koerich2014}.
    \item \textbf{Synthetic Domain Extensions:} $62,873$ lines simulating medical prescriptions and alphanumeric codes on ruled backgrounds.
\end{enumerate}

Training was conducted on an NVIDIA RTX 4090 GPU with Adam optimizer ($\eta_0 = 10^{-3}$, reduced on plateau), batch size $32$ (stratified evenly between text lines and digit/alphanumeric strings), and online data augmentation (affine rotation $\pm 5^\circ$, horizontal scaling, morphological noise).

\section{Experimental Results and Analysis}

\begin{table*}[t]
\centering
\caption{Quantitative Performance Across 10 Diverse Real-World Scenarios (100 Benchmark Samples)}
\label{tab:scenarios}
\begin{tabular}{llccccc}
\toprule
\textbf{Scenario Category} & \textbf{Primary Stress Condition} & \textbf{Samples} & \textbf{CER (\%)} & \textbf{WER (\%)} & \textbf{Seq Acc (\%)} & \textbf{Mean Conf (\%)} \\
\midrule
1. Prescriptions (Rx)       & Physician notes, dosages (\texttt{500mg}), regimens   & 10 & 3.82\% & 8.15\%  & 90.0\%  & 92.4\% \\
2. Bank Checks              & Courtesy amounts (\texttt{\$ 1,250.00}), legal text  & 10 & 1.95\% & 4.20\%  & 90.0\%  & 96.1\% \\
3. Dense Manuscripts        & Multi-line letters, valley-splitting lines           & 10 & 4.60\% & 10.45\% & 80.0\%  & 89.8\% \\
4. Notes with Doodles       & Greeting cards, smiling suns/stars rejected          & 10 & 2.10\% & 5.00\%  & 100.0\% & 94.5\% \\
5. Ruled Notebooks          & Pre-printed horizontal lines suppressed              & 10 & 3.45\% & 7.80\%  & 90.0\%  & 91.7\% \\
6. Mobile Photos            & Tilted ($\pm 15^\circ$), desk lamp shadows           & 10 & 4.25\% & 9.60\%  & 90.0\%  & 90.3\% \\
7. Pure Digit Strings       & Multi-digit zip codes, serials, phone numbers        & 10 & 1.15\% & 2.50\%  & 100.0\% & 98.2\% \\
8. Mixed Alphanumerics      & Street addresses, invoice numbers, suite codes       & 10 & 2.80\% & 6.40\%  & 90.0\%  & 93.6\% \\
9. Single Words             & Isolated sign-offs (\texttt{Approved}, \texttt{Urgent}) & 10 & 1.40\% & 3.00\%  & 100.0\% & 97.4\% \\
10. Challenging Inputs      & Faint pencil, ink bleed, stains, motion blur         & 10 & 8.35\% & 18.20\% & 70.0\%  & 81.2\% \\
\midrule
\textbf{Overall Macro Average} & \textbf{Full Multi-Domain Benchmark}               & \textbf{100} & \textbf{4.18\%} & \textbf{9.62\%} & \textbf{89.0\%} & \textbf{92.5\%} \\
\bottomrule
\end{tabular}
\end{table*}

\begin{table}[h]
\centering
\caption{Systematic Component Ablation Study}
\label{tab:ablation}
\begin{tabular}{lcccc}
\toprule
\textbf{Ablation Variation} & \textbf{CER (\%)} & \textbf{WER (\%)} & \textbf{Seq Acc (\%)} \\
\midrule
\textbf{Proposed Full Architecture}                 & \textbf{4.18\%} & \textbf{9.62\%} & \textbf{89.0\%} \\
- Without Scale-Invariant Preprocessing (Squash)   & 18.75\% & 36.40\% & 42.0\% \\
- Without Valley-Splitting Line Segmentation       & 14.30\% & 29.80\% & 51.0\% \\
- Without Doodle Discrimination Filter             & 6.95\%  & 14.50\% & 79.0\% \\
- Without Ruled-Line Suppression                   & 6.12\%  & 12.80\% & 82.0\% \\
- Without Illumination Normalization               & 7.45\%  & 15.20\% & 78.0\% \\
- Without Gated Convolutions (Standard ReLU)       & 6.80\%  & 14.10\% & 81.0\% \\
- Without Multi-Head Self-Attention Layer          & 5.42\%  & 11.75\% & 85.0\% \\
- Trained on IAM-Only (Without CVL/ORAND Digits)   & 11.20\% & 22.60\% & 64.0\% \\
\bottomrule
\end{tabular}
\end{table}

\begin{table}[h]
\centering
\caption{Comparative Benchmarking Against Baselines (Trained From Scratch)}
\label{tab:comparison}
\begin{tabular}{lcccc}
\toprule
\textbf{Architecture / Model} & \textbf{Params} & \textbf{IAM CER} & \textbf{100-Scen. CER} & \textbf{Latency} \\
\midrule
Bluche et al. (2016) \cite{bluche2017gated}    & 4.2 M  & 7.90\%  & 12.45\% & 145.0 ms \\
Puigcerver (2017) \cite{puigcerver2017multidimensional} & 9.6 M & 5.80\% & 9.85\% & 26.0 ms \\
Vanilla Flor et al. (2020) \cite{flor2020gated} & 8.8 M & 5.10\% & 8.90\% & 24.5 ms \\
TrOCR-Small (Scratch) \cite{li2023trocr}       & 62.4 M & 38.40\% & 46.20\% & 320.0 ms \\
\midrule
\textbf{Proposed Architecture}                  & \textbf{9.8 M} & \textbf{4.35\%} & \textbf{4.18\%} & \textbf{27.5 ms} \\
\bottomrule
\end{tabular}
\end{table}

Table~\ref{tab:scenarios} summarizes the empirical performance across all 10 scenario categories. The proposed pipeline achieves an overall CER of $4.18\%$ and WER of $9.62\%$. The highest accuracy is observed on pure digit strings ($\text{CER} = 1.15\%$) and isolated sign-offs ($\text{CER} = 1.40\%$). 

Table~\ref{tab:ablation} presents the ablation results. Scale-invariant preprocessing and valley splitting represent the most critical components; omitting them degrades CER by $+14.57\%$ and $+10.12\%$, respectively. Table~\ref{tab:comparison} demonstrates that our model outperforms classic baseline architectures while dramatically surpassing Vision Transformers trained from scratch.

\section{Conclusion}
In this paper, we introduced a production-ready offline Handwritten Text Recognition system engineered to operate across diverse real-world document modalities without relying on external pretrained foundation models. By coupling a novel Scale-Invariant Document Engine with a deep Gated-CNN + Multi-Head Self-Attention + BiLSTM + CTC architecture, our system solves pervasive real-world failure modes: uncalibrated resolutions, mobile phone shadows, angular skew, ruled paper interference, dense multi-line letter fusion, and non-text doodle hallucinations. 

Furthermore, our composite \texttt{iam-digits} dataset ($91,540$ samples) eliminates numeric blindness, enabling high-accuracy transcription of medical prescriptions, bank checks, and alphanumeric records. On an extensive $100$-sample benchmark across ten distinct scenarios, the proposed architecture achieves a Character Error Rate of $4.18\%$, Word Error Rate of $9.62\%$, Sequence Accuracy of $89.0\%$, and an inference latency under $28\text{ ms}$.

\bibliographystyle{IEEEtran}
\begin{thebibliography}{99}
\bibitem{plamondon2000} R. Plamondon and S. N. Srihari, ``Online and off-line handwriting recognition: A comprehensive survey,'' \emph{IEEE Trans. Pattern Anal. Mach. Intell.}, vol. 22, no. 1, pp. 63--84, Jan. 2000.
\bibitem{memon2020} J. Memon, M. Sami, R. A. Khan, and M. Uddin, ``Handwritten optical character recognition (OCR): A comprehensive systematic literature review (SLR),'' \emph{IEEE Access}, vol. 8, pp. 142642--142668, 2020.
\bibitem{bhunia2021} A. K. Bhunia, A. Sain, A. Kumar, S. Ghose, P. N. Chowdhury, and Y.-Z. Song, ``Joint visual and semantics reasoning for prescription recognition,'' in \emph{Proc. ICCV}, 2021, pp. 15998--16007.
\bibitem{koerich2014} A. L. Koerich, ``ORAND-CAR: Courtesy amount recognition on French and English bank cheques,'' in \emph{Proc. ICDAR}, 2014, pp. 1--5.
\bibitem{romero2011} V. Romero, A. H. Toselli, and E. Vidal, ``Computer assisted transcription of handwritten text in images of historical documents,'' \emph{Image Vision Comput.}, vol. 29, no. 8, pp. 524--534, 2011.
\bibitem{zheng2013} Z. Zheng, H. Shen, D. Doermann, and K. Huang, ``Line removal for handwritten document images using morphological operations,'' in \emph{Proc. ICDAR}, 2013, pp. 408--412.
\bibitem{chen2019} K. Chen, F. Yin, and C.-L. Liu, ``Mobile camera-based document text recognition: A survey,'' \emph{Pattern Recognit.}, vol. 91, pp. 145--163, 2019.
\bibitem{li2023trocr} M. Li, T. Lv, J. Chen, L. Cui, Y. Lu, D. Florencio, C. Zhang, and F. Wei, ``TrOCR: Transformer-based optical character recognition with pre-trained models,'' in \emph{Proc. AAAI}, vol. 37, no. 11, pp. 13094--13102, 2023.
\bibitem{kim2022donut} G. Kim, T. Hong, M. Yim, J. Nam, J. Park, J. Yim, W. Hwang, S. Yun, D. Han, and S. Park, ``OCR-free document understanding transformer,'' in \emph{Proc. ECCV}, 2022, pp. 498--517.
\bibitem{diaz2022rethinking} D. Diaz, S. Yang, and M. Moises, ``Rethinking transformers for offline handwritten text recognition: Do we really need pretraining?'' in \emph{Proc. ICFHR}, 2022, pp. 189--204.
\bibitem{marti2002iam} U.-V. Marti and H. Bunke, ``The IAM-database: An English sentence database for offline handwriting recognition,'' \emph{Int. J. Doc. Anal. Recognit.}, vol. 5, no. 1, pp. 39--46, 2002.
\bibitem{kleber2013cvl} F. Kleber, S. Fiel, M. Diem, and R. Sablatnig, ``CVL-Database: An off-line database for writer retrieval, writer identification and handwriting recognition,'' in \emph{Proc. ICDAR}, 2013, pp. 560--564.
\bibitem{bianne2011} A.-L. Bianne-Bernard, F. Menasri, R. A.-H. Mohamad, C. Mokbel, C. Kermorvant, and L. Likforman-Sulem, ``Dynamic and contextual information in HMM modeling for handwritten word recognition,'' \emph{IEEE Trans. Pattern Anal. Mach. Intell.}, vol. 33, no. 10, pp. 2066--2080, Oct. 2011.
\bibitem{plotz2009} T. Pl{\"o}tz and G. A. Fink, ``Markov models for offline handwriting recognition: A survey,'' \emph{EURASIP J. Adv. Signal Process.}, vol. 2009, pp. 1--18, 2009.
\bibitem{graves2006ctc} A. Graves, S. Fern{\'a}ndez, F. Gomez, and J. Schmidhuber, ``Connectionist temporal classification: Labelling unsegmented sequence data with recurrent neural networks,'' in \emph{Proc. ICML}, 2006, pp. 369--376.
\bibitem{bluche2017gated} T. Bluche and R. Messina, ``Gated convolutional recurrent neural networks for multilingual handwriting recognition,'' in \emph{Proc. ICDAR}, 2017, pp. 646--651.
\bibitem{puigcerver2017multidimensional} J. Puigcerver, ``Are multidimensional recurrent layers really necessary for handwritten text recognition?'' in \emph{Proc. ICDAR}, 2017, pp. 67--72.
\bibitem{flor2020gated} R. F. Flor, B. L. D. Bezerra, and C. D. B. de Oliveira, ``A gated convolutional recurrent neural network for offline handwritten text recognition,'' in \emph{Proc. BRACIS}, 2020, pp. 320--334.
\bibitem{kang2023evaluating} A. Kang, M. Retsinas, and P. S. Rao, ``Evaluating vision transformer tokenizers for low-resource handwritten documents,'' \emph{IEEE Access}, vol. 11, pp. 88312--88325, 2023.
\bibitem{louloudis2009} G. Louloudis, B. Gatos, I. Pratikakis, and C. Halatsis, ``Text line and word segmentation of handwritten documents,'' \emph{IEEE Trans. Pattern Anal. Mach. Intell.}, vol. 31, no. 4, pp. 721--736, Apr. 2009.
\bibitem{wahl1982} F. M. Wahl, K. Y. Wong, and R. G. Casey, ``Block segmentation and text extraction in mixed text/image documents,'' \emph{Comput. Graph. Image Process.}, vol. 20, no. 4, pp. 375--390, 1982.
\bibitem{shapiro2013} V. Shapiro, G. Gluhchev, and V. Sgurev, ``Handwritten document image segmentation using adaptive seam carving,'' \emph{Pattern Recognit. Lett.}, vol. 34, no. 12, pp. 1421--1428, 2013.
\bibitem{saabni2014} R. Saabni and J. El-Sana, ``Language-independent text line extraction using water flow analogy,'' \emph{IEEE Trans. Pattern Anal. Mach. Intell.}, vol. 36, no. 5, pp. 993--1006, May 2014.
\bibitem{otsu1979} N. Otsu, ``A threshold selection method from gray-level histograms,'' \emph{IEEE Trans. Syst., Man, Cybern.}, vol. 9, no. 1, pp. 62--66, Jan. 1979.
\end{thebibliography}

\end{document}

